% ---------------------------------------------------------------------------
% Results, part A: the model, the forward problem, the information budget.
% ---------------------------------------------------------------------------
\subsection{A seven-state model of the crypt, four disease regimes and a
ten-condition design}
\label{sec:res-model}

The model couples three loops (Fig~\ref{fig:design}A; Methods,
Eqs~\eqref{eq:b}--\eqref{eq:c}): a pro-stemness positive loop
\bcat{}$\to$\allowbreak MYC$\to$\allowbreak HOXA13$\to$\allowbreak\bcat{}; a differentiation loop in which RA
induces HOXA5, HOXA5 supports APC and APC degrades \bcat{}; and RA homeostasis
through its catabolising enzyme CYP26A1, which \bcat{} also induces. Disease
progression is a coupled sweep that raises WNT drive $W$ while lowering APC
functionality $\thP$ (Table~\ref{tab:regimes}); every other parameter is shared.
Along the sweep the \bcat{} level roughly triples, APC falls to a sixth of its
healthy level and the stemness index rises fourfold (Fig~\ref{fig:design}B). The
experiment that every inverse problem in this paper sees is the ten-condition
design of Fig~\ref{fig:design}C: six conditions vary only the retinoid input and
four perturb the WNT/MYC sub-network directly. All seven states are observed at
150 times per condition with Gaussian noise of standard deviation 0.002, 0.1--0.7\%
of the mean level of each state, except APC, whose level falls with WNT drive
(0.4\% in Normal to 2.6\% in Severe APC Loss).
This is a generous design, about $10^4$ observations per regime; how the
conclusions depend on it is examined in S1~Fig.

\begin{figure}[tbp]
\centering
\begin{minipage}[c]{0.58\linewidth}
{\sffamily\bfseries\large A}\\[-0.5ex]
\centering
\fig[width=0.97\linewidth]{fig01a_network.pdf}
\end{minipage}\hfill
\begin{minipage}[c]{0.39\linewidth}
\centering\footnotesize\setlength{\tabcolsep}{3pt}
\captionof{table}{\textbf{The four regimes.} Only $(W,\thP)$ change.}
\label{tab:regimes}
\begin{tabular}{lccl}
\toprule
regime & $W$ & $\thP$ & interpretation \\
\midrule
Normal & 0.80 & 1.00 & healthy crypt \\
Early Adenoma & 1.00 & 0.75 & WNT deregulation \\
Advanced Adenoma & 1.50 & 0.50 & established lesion \\
Severe APC Loss & 2.00 & 0.25 & carcinoma-like \\
\bottomrule
\end{tabular}
\end{minipage}\\[1.5ex]
\fig{fig01_design.pdf}
\caption{\textbf{Model, regimes and experimental design.} (A) The WNT--RA--HOX
network: one edge per term of Eqs~\eqref{eq:b}--\eqref{eq:c}, labelled with its
parameter; blue arrows activate, red bars repress. Orange badges mark the two
parameters that define the disease regimes; APC loss acts through the excess
degradation $\delta_P$. The dashed edge is the main learned term of the hybrid
models. (B) Reference trajectories (Radau, tolerance $10^{-10}$) of \bcat{}, APC
and the stemness index $S$ under the control protocol in the four regimes; grey
band, ATRA pulse. (C) The ten conditions of the base design: per condition, the
retinoid drive $\mu_R(t)$, the exogenous WNT input and the exogenous MYC input,
coloured where they differ from the control protocol (grey) and labelled with
their amplitudes.}
\label{fig:design}
\end{figure}

\subsection{The forward problem is not the main bottleneck}
\label{sec:res-forward}

Before asking how well a PINN estimates parameters, we ask whether it can
represent the dynamics at all, because an inverse failure could simply be a
representation failure. With the parameters known and 40 noise-free observations
per state, a Fourier-feature PINN reproduces all seven states in all four regimes
(Fig~\ref{fig:forward}A) with a mean relative $L^2$ error of 2.83\%, 1.76\%,
2.09\% and 2.95\% (grand mean 2.41\%, or 1.59\% without APC), against 1.06\% for
the same network trained on 100 evenly spaced labels (Fig~\ref{fig:forward}B).
The largest relative error is in APC in the high-WNT regimes, where APC settles
near zero and the denominator of the relative error is small (13.5\% relative,
1.3\% normalised RMSE). The network is expressive enough for the trajectories,
including the circadian oscillation and the treatment pulse. Its error is
larger than the 0.1--2.6\% observation noise of the inverse problem, so
representation is not free, but it is a few percent where the parameter errors
of the inverse PINN below are tens of percent; the representation of the states
is not what limits them (Section~\ref{sec:res-sgn} returns to this).

\begin{figure}[tbp]
\centering
\fig{fig02_forward.pdf}
\caption{\textbf{A PINN represents the dynamics.} (A) Forward PINN (dashed,
regime colour) trained on 40 noise-free observations (open circles), the initial
condition and the ODE residual, against the Radau reference (black), for the
seven states in the Normal and Severe APC Loss regimes. (B) Relative $L^2$
error per state and regime of this sparse-data PINN (filled) and of the same
network trained on 100 evenly spaced labels (open). Mean errors 2.83/1.76/2.09/2.95\%
(grand mean 2.41\%, against 1.06\%); the APC error in the high-WNT regimes is
large in relative terms because APC approaches zero there.}
\label{fig:forward}
\end{figure}

\subsection{The information budget of the design}
\label{sec:res-budget}

The Fisher information of the ten-condition design at the true parameters,
Eq~\eqref{eq:fim}, is a single Jacobian computation that needs no estimator.
In every regime its 35 informative eigenvalues span 5.7 to 8.3 orders of
magnitude, and one further eigenvalue sits 7.8 to 10.9 orders below them, at the
level of finite-difference round-off: a structural null
direction (Fig~\ref{fig:info}A). The null direction is exactly
the analytic degeneracy of Eq~\eqref{eq:p}: it lies in the $(\log\thP,\log\dP)$
plane to within $10^{-6}$, with loadings that match the tangent of the level
set $\dP(1-\thP)=\text{const}$ (Fig~\ref{fig:info}B). Its $\thP$ loading rises
from 0 to 0.95 across the regimes, but that is algebra, the ratio
$(1-\thP):\thP$, not a sign of saturation; we mention it because it is easily
mistaken for one. The Fisher matrix finds the one structural fact about the
model without being told, which we take as a check of the computation.

With $\dP$ known, the CRLB of the remaining 35 parameters is the information
budget of the design (Fig~\ref{fig:info}C). In the two low-WNT regimes every
parameter has a relative standard deviation below 1.3\%; the budget then
shrinks with WNT drive, to 31 of 35 parameters below 5\% in Advanced Adenoma
and 25 of 35 in Severe APC Loss, where $\thP$ (57\%), $\rho_{13}$ (21\%),
$\varepsilon_P$ (18\%) and $a_C$ (16\%) are the least informed. Translated into
the currency the literature on this model has used, an efficient estimator would
recover on average 34.0, 34.0, 33.4 and 31.1 of the 34 parameters outside the
degenerate pair to within 10\% (the CRLB-expected count $\mathbb{E}[\Ncount]$,
Methods), 132.5 of 136 in total (Fig~\ref{fig:info}D). A naive reading of recovery counts, of the kind the PINN results below produce,
would conclude that the data support a dozen or so parameters per regime. The
bound says the data support nearly all of them.

\begin{figure}[tbp]
\centering
\fig{fig03_information.pdf}
\caption{\textbf{The information budget of the ten-condition design.} Fisher
information of the 36 log-parameters at the truth (150 observations per
condition, seven states, $\sigma=0.002$). (A) Eigenvalue spectra normalised by
the largest eigenvalue: one structural null direction in every regime. (B)
Loadings of the null eigenvector on $\log\thP$ (light) and $\log\dP$ (dark);
crosses, the tangent of the level set $\dP(1-\thP)=\text{const}$, which they
match exactly. (C) Cram\'er--Rao relative standard deviation of each parameter
with $\dP$ known, sorted by its value in Severe APC Loss; background bands,
identifiable ($<5\%$), weak (5--50\%), non-identifiable ($>50\%$). (D)
CRLB-expected number of the 34 parameters outside the degenerate pair that an
efficient estimator recovers within 10\% (bars), and the number with CRLB
standard deviation below 5\% (ticks).}
\label{fig:info}
\end{figure}
